\chapter{Haemoptysis (Coughing Up Blood)}

\section*{Definition}
Expectoration of blood or blood-tinged sputum from the respiratory tract below the larynx, distinguishing it from hematemesis (gastrointestinal) and epistaxis (nasal).

\section*{When to See a Doctor}
See your doctor if:
\begin{itemize}
  \item Any episode of hemoptysis requires medical evaluation; large volume (>100mL) or hemodynamic instability requires emergency care
\end{itemize}

\section*{How It Develops}
Bleeding originates from bronchial arteries (systemic circulation, 90% of significant hemoptysis) or pulmonary arteries (low pressure, less common but catastrophic). Causes include airway inflammation (bronchitis, bronchiectasis), parenchymal disease (pneumonia, TB, malignancy), vascular disorders (AVM, PE), or bleeding diatheses.

\section*{Causes}
\begin{itemize}
  \item Acute bronchitis (most common cause)
  \item Bronchiectasis
  \item Tuberculosis
  \item Lung cancer or endobronchial tumor
  \item Pneumonia (including lung abscess)
  \item Pulmonary embolism with infarction
  \item Aspergilloma (fungus ball in pre-existing cavity)
  \item Bronchial adenoma or carcinoid tumor
  \item Pulmonary arteriovenous malformation
  \item Autoimmune disease (Wegener granulomatosis, Goodpasture syndrome)
\end{itemize}

\section*{Related Symptoms}
\begin{itemize}
  \item Cough
  \item Shortness of breath
  \item Chest pain
  \item Fever
  \item Weight loss
\end{itemize}

\section*{Medical Evaluation}
\begin{itemize}
  \item Your doctor will take a history of the amount, frequency, and triggers of the bleeding.
  \item Chest X-ray is the initial imaging study to evaluate for pneumonia, mass, or cavity.
  \item CT chest with contrast provides detailed evaluation of lung parenchyma and bronchi.
  \item Sputum culture and cytology help identify infection or malignant cells.
  \item Bronchoscopy may be performed for direct visualization and biopsy of airway lesions.
\end{itemize}

\section*{Treatment Options}
\begin{itemize}
  \item Antibiotics for infectious causes (bronchitis, pneumonia, tuberculosis).
  \item Bronchial artery embolization for life-threatening hemoptysis.
  \item Surgical resection (lobectomy, segmentectomy) for lung cancer or bronchiectasis.
  \item Antifungal therapy for aspergilloma.
  \item Immunosuppressive therapy for autoimmune causes (vasculitis, Goodpasture syndrome).
\end{itemize}

\section*{Self-Care and Home Management}
\begin{itemize}
  \item Even a small amount of coughing up blood warrants medical advice; do not assume it is benign or due to bronchitis.
  \item Sip cool water to soothe the throat and help clear blood-tinged mucus.
  \item Avoid coughing forcefully — gentle coughing is sufficient to clear the airways.
  \item Do not take blood-thinning medications (aspirin, NSAIDs) without consulting your doctor.
  \item Seek emergency care for large volume hemoptysis (more than a few teaspoons) or difficulty breathing.
\end{itemize}

\section*{References}

\begin{thebibliography}{99}
\bibitem{haemoptysis:harrison-hemopt} Longo DL, et al. \emph{Harrison's Principles of Internal Medicine}. 21st ed. McGraw-Hill; 2022. Chapter 49: Hemoptysis.
\bibitem{haemoptysis:uptodate-hemopt} Earwood JS, et al. Evaluation of hemoptysis in adults. In: UpToDate, Post TW (Ed), UpToDate, Waltham, MA; 2025.
\bibitem{haemoptysis:lordan} Lordan JL, et al. The management of hemoptysis. \emph{Eur Respir Rev}. 2023;32(167):220179.
\bibitem{haemoptysis:sakr} Sakr L, et al. Massive hemoptysis: diagnosis and treatment. \emph{Chest}. 2022;161(3):673-684.
\bibitem{haemoptysis:fishman} Grippi MA, et al. \emph{Fishman's Pulmonary Diseases and Disorders}. 6th ed. McGraw-Hill; 2023. Chapter 57: Hemoptysis.
\bibitem{haemoptysis:bnf-tb} World Health Organization. \emph{WHO Guidelines on Tuberculosis}. WHO; 2024. Chapter 4: Diagnosis of Pulmonary TB.
\end{thebibliography}
